%% ==============================================================
%% КОНФИГУРАЦИЯ — Вариант 10
%% Геометрия A, тип 10 (Трапецеидальный (прямоугольный))
%% Сгенерировано автоматически — не редактировать вручную
%% ==============================================================

% --- Информация о студенте ---
\newcommand{\varNumber}{10}
\newcommand{\varStudent}{Михайлов~Д.А.}
\newcommand{\varGroup}{ЗНБ21-02Б}
\newcommand{\varSupervisor}{Башмур~К.А.}

\newcommand{\varStudentID}{082159252}

% --- Геометрия подшипника (геометрия A) ---
\newcommand{\varR}{0.035}
\newcommand{\varRmm}{35}
\newcommand{\varC}{5e-05}
\newcommand{\varCmkm}{50}
\newcommand{\varL}{0.056}
\newcommand{\varLmm}{56}
\newcommand{\varLD}{0.8}
\newcommand{\varGeomLetter}{A}

% --- Тип микрорельефа (тип 10) ---
\newcommand{\varDimpleType}{10}
\newcommand{\varDimpleTypeName}{Трапецеидальный (прямоугольный)}
\newcommand{\varDimpleTypeGenitive}{трапецеидальным}

% --- Параметры микрорельефа ---
\newcommand{\varDimpleA}{0.00238}
\newcommand{\varDimpleAmm}{2.38}
\newcommand{\varDimpleB}{0.00218}
\newcommand{\varDimpleBmm}{2.18}
\newcommand{\varDimpleHp}{1e-05}
\newcommand{\varDimpleHpmkm}{10}
\newcommand{\varDimpleHpRatio}{0.2}

% --- Формула зазора ---
\newcommand{\varDimpleFormula}{
    \Delta h = h_p \text{ при } r_\infty\leq0.6;\; h_p(1-r_\infty)/0.4 \text{ при } 0.6<r_\infty\leq1
}

% --- Общие параметры ---
\newcommand{\varN}{2980}
\newcommand{\varEta}{0.01105}
\newcommand{\varNphi}{8}
\newcommand{\varNZ}{11}
\newcommand{\varPhiRange}{$90^\circ\text{ - }270^\circ$}
